% --- INICIO DE PLANTILLA PERSONALIZADA ---
\documentclass[oneside,11pt]{Latex/Classes/PhDthesisPSnPDF}

% --- Paquetes Originales ---
\usepackage{blindtext}
\usepackage{actuarialangle} 
\usepackage{tabularx}
\usepackage{float}
\usepackage{multirow, array}
\usepackage[usenames,dvipsnames,svgnames,table]{xcolor}
\usepackage{longtable}
\usepackage{latexsym,amsmath,amssymb,amsfonts}
\usepackage{enumitem}
\usepackage{xparse}
\usepackage{booktabs}
\usepackage[round, sort, numbers]{natbib}  
\usepackage{adjustbox}
\usepackage[utf8]{inputenc}
\usepackage{caption}
\usepackage{graphicx}
\usepackage{hyperref}
\usepackage{cleveref}

% Definición de colores
\definecolor{miblue}{RGB}{68, 114, 196}

% --- CAMBIO CRÍTICO 1: Usar \input para macros en el preámbulo ---
% \include da problemas antes del begin document. \input es seguro.
\input{Latex/Macros/MacroFile1}

% --- Definiciones de seguridad (por si la clase falla al cargar alguna variable) ---
\providecommand{\facultad}[1]{\def\lafacultad{#1}}
\providecommand{\escudofacultad}[1]{\def\elescudo{#1}}
\providecommand{\degree}[1]{\def\elgrado{#1}}
\providecommand{\director}[1]{\def\eldirector{#1}}
\providecommand{\lugar}[1]{\def\ellugar{#1}}
\providecommand{\degreedate}[1]{\def\lafecha{#1}}

% --- Configuración para bloques de código de R (Pandoc) ---

%%%%%%%%%%%%%%%%%%%%%%%%%%%%%%%%%%%%%%%%%%%%%%%%%%%%%%%%%%%%%%%%%%%%%%%%%%%%%%%%
%                         DATOS (Conectados a RMarkdown)                       %
%%%%%%%%%%%%%%%%%%%%%%%%%%%%%%%%%%%%%%%%%%%%%%%%%%%%%%%%%%%%%%%%%%%%%%%%%%%%%%%%
\title{Análisis Descriptivo de los Factores Determinantes en los
Cultivos}
\author{Joshua Santeliz \\ Miguel Urbina \\ }

% Lógica para Facultad: Si la defines en el YAML la usa, si no, usa la default

\facultad{Facultad de Ciencias Económicas y Sociales \\ Escuela de
Estadistica y Ciencias Actuariales}
  \escudofacultad{Latex/Classes/Escudos/faces}

\degree{Computación I}
\director{Jesus Ochoa \\ Oliver Triveño}
\degreedate{Febrero 2026} 
\lugar{Caracas}

\portadatrue 

% Metadatos PDF
\keywords{}
\subject{}

% --- CORRECCIÓN PARA PANDOC NUEVO (Imágenes) ---
\makeatletter
\providecommand{\pandocbounded}[1]{#1}
\makeatother


%%%%%%%%%%%%%%%%%%%%%%%%%%%%%%%%%%%%%%%%%%%%%%%%%%%%%
%                   DOCUMENTO                       %
%%%%%%%%%%%%%%%%%%%%%%%%%%%%%%%%%%%%%%%%%%%%%%%%%%%%%
\begin{document}

\maketitle

%%%%%%%%%%%%%%%%%%%%%%%%%%%%%%%%%%%%%%%%%%%%%%%%%%%%%
%                  PRÓLOGO                          %
%%%%%%%%%%%%%%%%%%%%%%%%%%%%%%%%%%%%%%%%%%%%%%%%%%%%%
\frontmatter

% Puedes agregar dedicatorias desde el YAML si quieres, o dejarlas fijas aquí

%%%%%%%%%%%%%%%%%%%%%%%%%%%%%%%%%%%%%%%%%%%%%%%%%%%%%
%                   ÍNDICES                         %
%%%%%%%%%%%%%%%%%%%%%%%%%%%%%%%%%%%%%%%%%%%%%%%%%%%%%
\setcounter{secnumdepth}{3}
\setcounter{tocdepth}{3}

\tableofcontents



\cleardoublepage

  \cleardoublepage       % Empieza en página derecha (estándar en tesis)
  \chapter*{Resumen}     % Crea el título "Resumen" sin numeración (Capítulo *)
  \addcontentsline{toc}{chapter}{Resumen} % (Opcional) Lo añade al índice
  
  Resumen del trabajo de
investigación\ldots{}             % Aquí se pega el texto que escribiste en el YAML
  
  \cleardoublepage

%%%%%%%%%%%%%%%%%%%%%%%%%%%%%%%%%%%%%%%%%%%%%%%%%%%%%
%                   CONTENIDO (El Corazón)          %
%%%%%%%%%%%%%%%%%%%%%%%%%%%%%%%%%%%%%%%%%%%%%%%%%%%%%
\mainmatter
\def\baselinestretch{1.5}

% --- CAMBIO CRÍTICO 2: Aquí RMarkdown inyecta todo tu texto ---
\chapter*{Introducción}

Introducción del trabajo de investigación

\chapter{Planteamiento del Problema}

En el mundo agrícola hay muchos temas por tratar, desde la salud
nutricional hasta la gastronomía industrial, siendo temas de vital
importancia para la vida humana. En este estudio analizaremos
específicamente a la agricultura de la India en el ámbito del cultivo,
resaltando que un manejo eficiente no solo garantiza cosechas óptimas,
sino que genera seguridad alimentaria y eleva la calidad de vida de las
personas. Los cultivos son influenciados por factores como el clima, los
ciclos agrícolas, fertilizantes y pesticida, estos factores son
importantes ya que a las empresas les interesas esas influencias para
mejorar su gestión de inversión, bajando la incertidumbre al describir
estas influencias.\\

Por lo antes mencionado surgen problemas por ejemplo en la parte del
clima como fenómenos naturales, problemas de contaminación que podría
afectar ya sea al clima o a los mismos cultivos, estos problemas se
resuelven luego de modelar esta investigación donde podremos comprender
mejor la situación y se recomendará posibles decisiones a tomar. Debido
a esto, se plantea las siguientes preguntas de investigación:\\

\textbf{1.} ¿De qué manera influye el clima en la calidad de los
cultivos?\\

\textbf{2.} ¿Por qué son necesario los fertilizantes y los pesticidas en
una temporada de cultivo? ¿Afectaría la salud de la población de la
India a nivel nutricional?\\

\textbf{3.} ¿Cuáles son los cultivos que más se dan según las temporadas
del año?\\

\section{Justificación}

La justificación de este trabajo de investigación reside en la necesidad
de aplicar técnicas estadísticas rigurosas y herramientas de software
como R y Python para mitigar la incertidumbre de sembrar los cultivos en
el sector agrícola de la India. Desde el punto de vista práctico, los
resultados permitirán a las Industrias tomar decisiones basadas en datos
(Data-Driven Decisions) sobre que materiales invertir o en qué temporada
del año se debe desarrollar. Con una metodología replicable utilizando R
y LaTeX para el análisis exploratorio y descriptivo de grandes volúmenes
de datos.\\

\section{Objetivos}

\subsection{Objetivo General}
\begin{itemize}
  \item{Realizar un análisis estadístico descriptivo de datos de cultivos en diferentes temporadas de la India durante el período 2010-2020.  para identificar y describir las características más influyentes en los cultivos.}
\end{itemize}

\subsection{Objetivos Específicos}
\begin{itemize}
  \item{1.  Analizar la variabilidad pluviométrica estacional mediante el uso de indicadores estadísticos, con el fin de determinar una óptima gestión del recurso hídrico en cada ciclo de cultivo}
  \item{2.  Observar la distribución de los factores que afectan a los cultivos por temporada del año.}
  \item{3.  Categorizar temporadas de mayor y menor demanda tanto de pesticidas como de fertilizantes}
\end{itemize}

\section{Cobertura}

\subsection{Horizontal}

La cobertura horizontal abarca las características de cada cultivo
registrado, donde se consideran las dimensiones descriptivas del cultivo
y las métricas de rendimiento agrícola.\\

\begin{table}[ht]
\centering
\caption{Variables consideradas} 
\resizebox{10cm}{!} {
\begin{tabular}{ccc}
  \hline
\ \ \ \ \ NOMBRE DE LA VARIABLE \\ 
  \hline
   VARIABLE 1 \\ 
VARIABLE 2\\ 
VARIABLE n\\ 
   \hline
\end{tabular}
}
\end{table}

\subsection{Vertical}

Describir la cobertura Vertical del trabajo de investigación

\section{Periodo de Referencia}

Describir la fuente y el Periodo de Referencia utilizado en el trabajo
de investigación

\section{Variables}

Describir la variable en estudio en el trabajo de investigación

\chapter{Marco Teórico}

\section{Bases Teóricas}

Resumen de las bases teoricas explicadas en el capitulo y su
importancia\\

\subsection{Definición 1}

Descripción de la definición 1\\

\subsection{Definición 2}

Descripción de la definición 2\\

\subsection{Definición n}

Descripción de la definición n\\

\chapter{Marco Metódico}

En esta sección se presentan las metodologías y/o test necesarios
relacionados con la práctica divulgada en el marco teórico, así como las
fuentes de datos.

\section{Bases metódicas utilizadas en el trabajo de investigación}

Desarrollo\\

\subsection{Item 1}

Descripción\\

\subsection{Item 2}

Descripción\\

\subsubsection{Item 2.2}

Descripción\\

\subsection{Item n}

Descripción\\

\chapter{Análisis de Resultados}

\{\small Este capítulo presenta los resultados obtenidos al aplicar las
metodologías descritas en el capítulo anterior, así como del análisis
estadístico descriptivo de las variables estudiadas y las ventajas y
desventajas al aplicar dichos tests.\}

\section{Presentación de resultados 1}

Descripción\\

\section{Presentación de resultados 2}

Descripción\\

\section{Presentación de resultados n}

Descripción\\

\chapter{Conclusiones y Recomendaciones}

\textbf{1.-} Conclusion 1 del trabajo de investigación.\\

\textbf{Recomendación} Recomendación 1 del trabajo de investigación.\\

\textbf{2.-} Conclusion 2 del trabajo de investigación.\\

\textbf{Recomendación} Recomendación 2 del trabajo de investigación.\\

\textbf{n.-} Conclusion n del trabajo de investigación.\\

\textbf{Recomendación} Recomendación n del trabajo de investigación.\\

\appendix
\renewcommand{\thechapter}{\Alph{chapter}}

\chapter{Anexo A}

\chapter{Anexo B}

\chapter{Anexo C}

%%%%%%%%%%%%%%%%%%%%%%%%%%%%%%%%%%%%%%%%%%%%%%%%%%%%%
%                   REFERENCIAS                     %
%%%%%%%%%%%%%%%%%%%%%%%%%%%%%%%%%%%%%%%%%%%%%%%%%%%%%
% Mantenemos la estructura natbib de tu plantilla
\renewcommand{\bibname}{Lista de Referencias}
\bibliographystyle{apalike} 
\bibliography{bibliografia.bib} 

%%%%%%%%%%%%%%%%%%%%%%%%%%%%%%%%%%%%%%%%%%%%%%%%%%%%%
%                   APÉNDICES                       %
%%%%%%%%%%%%%%%%%%%%%%%%%%%%%%%%%%%%%%%%%%%%%%%%%%%%%

\end{document}
